% !TEX program = xelatex
% 由 tactic_translate.export_plays_pdf 產生；勿手動編輯
% 編譯：cd final_document && xelatex A001-A006_戰術語言.tex
\documentclass[10pt]{article}
\usepackage[a4paper,margin=1.6cm]{geometry}
\usepackage{fontspec}
\usepackage{xeCJK}
\usepackage{titlesec}

\setmainfont{Times New Roman}
\IfFontExistsTF{Songti TC}{
  \setCJKmainfont{Songti TC}
  \setCJKsansfont{Heiti TC}
}{
  \setCJKmainfont{FandolSong}
  \setCJKsansfont{FandolHei}
}

\titlespacing*{\subsection}{0pt}{0.6em}{0.25em}
\titleformat{\subsection}{\normalsize\bfseries}{}{0em}{}
\setlength{\parskip}{0.25em}
\pagestyle{empty}

\begin{document}

{\Large\bfseries 戰術復刻　A001–A006 戰術語言\par}
\vspace{0.6em}

\subsection*{A001　題型 A001：邊路經典爆破}
\noindent\textbf{戰術語言}　邊後衛（FB）在右邊路（後段）接球。邊後衛傳球給前腰（AM），前腰在右肋接球。前腰傳球給邊後衛（FB），邊後衛在右路底線接球。邊後衛傳中給前鋒（ST），前鋒在禁區中央得分。
\vspace{0.5em}

\subsection*{A002　題型 A002：大範圍撕裂與倒三角}
\noindent\textbf{戰術語言}　中場（CM）在右路中路接球。中場強弱邊轉移給左邊鋒（LW），左邊鋒在左路底線接球。左邊鋒倒三角傳球給前鋒（ST），前鋒在弧頂前緣接球。前鋒傳球給跑進禁區的左邊鋒（LW），左邊鋒在禁區中央得分。
\vspace{0.5em}

\subsection*{A003　題型 A003：中路三人過度}
\noindent\textbf{戰術語言}　後腰（DM）在後場禁區接球。後腰從己方禁區直塞給前鋒（ST1），前鋒在弧頂前緣接球。前鋒回做球給前腰（AM），前腰在左肋深處接球。前腰直塞給前鋒（ST2），前鋒在禁區中央射正。
\vspace{0.5em}

\subsection*{A004　題型 A004：第三人破高位壓迫}
\noindent\textbf{戰術語言}　中後衛（CB）在後場禁區接球（遭高位壓迫）。中後衛傳球給後腰（DM），後腰在後場中路接球。後腰傳球給邊後衛（FB），邊後衛在右邊路（後段）接球（第三人出球脫困）。邊後衛直塞給前腰（AM），前腰在禁區中央射正。
\vspace{0.5em}

\subsection*{A005　題型 A005：RW／FB／CM 倒腳推進射門}
\noindent\textbf{戰術語言}　右邊鋒（RW）在右邊路（中段）接球。右邊鋒傳球給邊後衛（FB），邊後衛在後場右角接球。邊後衛傳球給中場（CM），中場在後場中路接球。中場傳球給右邊鋒（RW），右邊鋒在右路中路接球。右邊鋒傳球給中場（CM），中場在弧頂前緣接球。中場回做球給邊後衛（FB），邊後衛在後場中路接球。邊後衛傳球給中場（CM），中場在右邊路接球。中場傳球給右邊鋒（RW），右邊鋒在弧頂前緣射正。
\vspace{0.5em}

\subsection*{A006　題型 A006：邊路起腳射門（左路）}
\noindent\textbf{戰術語言}　中場（CM）在後場中路接球。中場傳球給前腰（AM），前腰在左肋接球。前腰傳球給左邊鋒（LW），左邊鋒在左邊路（中段）接球。左邊鋒傳球給前腰（AM），前腰在弧頂前緣接球。前腰傳球給左邊鋒（LW），左邊鋒在左邊路（前段）射正。
\vspace{0.5em}

\end{document}
